\chapter{Arrangementer}

På \TKET elsker vi at hygge os i godt selskab. Derfor tøver vi ikke med at opfinde undskyldninger for at mødes ved enhver given lejlighed. Nogle af disse undskyldninger er med tiden blevet til faste arrangementer hvis de har været en success to gang i streg inden for en fornuftig periode (typisk værende to år eller hvor lang tid det nu kræver at opleve succesen to gange).

I dette afsnit gennemgås de fleste af arrangementerne, som har været arrangeret i \TK-regi -- både nuværende og tidligere. Første afsnit beskriver de nuværende arrangementer opsat efter deres omtrentlige tidslige placering i løbet af bestyrelsesåret, mens andet afsnit beskriver tidligere afholdte arrangementer.

%%%%%
\section{Nuværende arrangementer}

\subsection{Torsdagsmøde}
Bestyrelsesmøde afholdes hver torsdag klokken halv otte -- for det meste kl. 19.30, men enkelte gange kl. 7.30. Her er alle velkomne. \TKET giver alle fremmødte en punkt-0-drikkevare. Husk at høre den!

\subsection{HSTR}
Hemmelig Skovtur Til Rønde, også tidligere kendt som Løvfaldstur. HSTR er en løvfaldstur til en hemmelig destination (Rønde), hvor den nye BESTyrelse med hjælp fra pegeFUer forsøger at holde styr på en flok fulde hængere på tur gennem en skov, hvor der bliver serveret øl og kage, er mulighed for at bade, og slutter af på en kro, hvor der serveres biksemad. En sjov tradition er her FUernes kagespisningskonkurrence. Den startede i 2011, hvor den ``forsvarende'' mester, daværende TOFUAB, udfordrede de nyvalgte FUer i den ``traditionsrige FUkagespisningskonkurrence''. Siden da er ``traditionen'' faktisk blevet en tradition, dog har konkurrencen ændret mængden af kage, som skal spises, samt om der spises smørstang eller dagmartærte.

\subsection{Reception}
To gange årligt holder \TKET åbent hus, hvor ansatte, studerende og andre VIP/TAP'er inviteres til smørrebrød, portvin og en hyggelig eftermiddag på \TKET.

\subsection{Fødselsdagstorsdagsmøde}
\TKET har fødselsdag d. 11. oktober og det skal fejres! Derfor afholdes der fødsesdagstorsdagsmøde en torsdag omkring dette tidspunkt. I god fødselsdagsstil serveres der her kagemand samt rød sodavand, og der synges fødselsdagssang for \TKET.

\subsection{Halloweenfest}
Første fest med den nye bestyrelse. I 25/26 blev Halloweenfesten indført som en kombination mellem Stiftelsesfesten og Karnevalsfesten. Festen inkluderer mulighed for spisning, hvor der undervejs foregår spas og løjer fra BEST side, og efterfølgende fest med kostumekonkurrence.

\subsection{Foredrag}
Hvert kvartal i bestyrelsesåret byder på et foredrag, hvor de studerende kan komme og blive underholdt med et emne der ikke er fagligt relevant. Foredragene spænder vidt i emner og fra de sidste ti år har der blandt andre være foredrag ved følgende:

\begin{multicols}{3}
  \begin{itemize}
    \item Søren Ryge
    \item Sigurd Barrett
    \item Sebastian Klein
    \item Ulrich Larsen
    \item Vicky Knudsen
    \item Nikolaj Sonne
    \item Miss Megan Moore
    \item Nikolaj Kirk
    \item Per Vers
    \item Thomas Skov
  \end{itemize}
\end{multicols}

Der afholdes Café Bagefter bagefter.

\subsection{Café Bagefter}
Efter hvert foredrag inviterer \TKET både studerende og foredragsholder til at nyde en drikkevare og godt selskab i Café Bagefter, hvor der er mulighed for at spille brætspil og diskutere det netop afsluttede foredrag.

\subsection{BODYCRASHING}
Bræk! Og Desuden Yder \CERM Rationel Afskaffelse Af Sultne Homosapiens I Naturvidenskabelige Gemakker. BODYCRASHING er et berømt og berygtet arrangement, hvor folk konkurrerer mere eller mindre frivilligt i kunsten at indtage mest mulig mad og drikke på mindst mulige tid. Det hele startede i 1985 hvor daværende FORM og NF besluttede sig for at stjæle idéen fra Datalogisk Fredagsbar, og gå over til netop den fredagsbar og bunde fløde som happening for \TKETs rafleturnering. Til turneringen udviklede duellen sig til en konkurrence kaldet ``Den Dræbende Fløde'' -- først bundes $\nicefrac{1}{2}$ L føde efterfulgt af en lunken top og sidst igen $\nicefrac{1}{2}$ L fløde. Vinderen er den hurtigste, hvorved det kunne etables hvem der havde magten i \TKET, FORM eller NF. Rafleturneringen var et kæmpe hit, så fra senere år kom BODYCRASHING på programmet som et arrangement for sig selv med flere forskellige discipliner. Her dystede sultne studerende i at indtage store mængder mad, ulækre kombinationer af mad, eller begge dele, så hurtigt som muligt.\par
Af populære discipliner kan nævnes:

\begin{itemize}
  \item Cocio på tid: Drik så mange 24 cl flasker Cocio man kan på tre minutter. I ti år holdt ``Cocio-Kenneth'' rekorden med 29 styk, men denne blev slået i 2023, da formen på flasken ændrede sig, og tillod INKA 24/25 at sætte ny rekord på 33 flasker. Samme mand holder i dag rekorden på 34 styk, slået i 2025.
  \item $2^4-3$: Drik 2 dl af følgende i rækkefølge: Smeltet jordbæris, olie/eddike-dressing, svineboullion, svagt pisket fløde, frisk cola, doven øl, sildelage, soyamælk, grapejuice, Cocio, kold automatte, drikkeyoghurt og tomatjuice. Samtidig lytter man til ``Kur mod tømmermænd'' fra Den Kummerlige Trio. Gruppen skulle have fået idéen til sangen fra en håndfuld hængere på Smukfest. Rekorden er  35 sekunder, sat af Jens i 2016.
  \item Blendet grillbar: En bearnaise-burger, en gang pomfritter, en ristet hotdog med det hele og én Cocio blendet sammen, blødet op med cola og serveret i en glasstøvle. Rekorden er på 24,24 sekunder, sat af Mathias i 2017.
  \item Akademiker ville vælge Whiskas: Kattemad og kærnemælk. Rekorden er på 43,2 sekunder, sat af Jens i 2017.
\end{itemize}

Rammerne omkring arrangementet er vokset. Disciplinerne foregår på \CERM{}s scene, dele af lokalet er spærret af til hurtig udgang til deltagerne, samt den ildevarslende ``splash-zone''. FUer iklædes hermetisk lukkede dragter og stiller sig solidariske til rådighed til at fange opkast i hvide affaldsspande. Medicinstuderende lokkes med en kasse guldøl til at supplere ekspertviden, så ingen deltager udsættes for noget, der i sandhed er sundhedsskadeligt

\subsection{J-$\nicefrac{1}{2}$-dag}
Den første fredag i november er J-dag, hvor juleøllen frigives. I den anledning afholdes J-$\nicefrac{1}{2}$-dag i Matematisk Kantine omkring middagstid, hvor der synes sange og drikkes juleøl fra det svundne år.

\subsection{Juleklip}
\TKET skal pyntes op til jul. Af denne grund arrangerer PR juleklip, hvor man kan sidde og klippe julehjerter og andet julepynt eller folde fine julestjerner. Som tak for fremstilling af et stykke pynt til ophængning på \TKET sørger PR for rigeligt med æbleskiver og glögg.

\subsection{Julerevy og -fest}
Den sidste studiedag i vintersemesteret afholder \TKET en Revy. Revyen indeholder sketches, sange og (pause-)fisk som kan hænde at være baseret på hændelser der er sket på både nationalt og internationalt plan i løbet af det sidste halve år. Udover sketches vil der til denne revy også forekomme opvisning af cabaret dans, ofte mere behåret end man forventer. Revyen afsluttes med en hvid x-mas og efterfølgende fest.

\subsection{Luciaoptog}
En undskyldning for at gå i kjole og synge falsk. Foregår normalt i december på (eller lige omkring) luciadag. Den siddende FUAN (eller en anden mandig mand) er luciabrud. Luciaoptoget snoer sig fra Matematisk institut op forbi Fysik og Nano og ender i Studieadministrationens bygninger til meget underholdning for både studerende og ansatte på universitetet.

\subsection{VKF}
Verdens Kedeligste Foredrag. Tre respekterede forelæsere udfordres til ikke at sige noget spændende i 20 minutter. FU'erne står for arrangementet, og de konkurrerer desuden om, hvem der kan lave den bedste plakat. BEST bager kage.

\subsection{Verdens Mindste Tour de Fredragsbar (VMTdF)}
Verdens mindste Tour de Fredagsbar er et samarbejde med fredagsbarer/festforeninger for studierne under \TKET, heriblandt Kalkulerbar, Fysisk Fredagsbar, Fredagscaféen, Nanorama og MØF.

Alle har hver deres bod, hvor de studerende nemt kan prøve nogle af deres specialiteter, om det er specialøl fra Datbar og FFB, shots fra Kalkulerbar og Nanorama, eller hvad de ellers finder på. Hver bar/forening holder også et event i løbet af aftenen.

\subsection{47. Dag}
Den 47. Dag er en historiefyldt tradition, som har været fejret ca. altid. Afhængigt af om et givet år er et skudår falder d. 47 dag på den 16. februar, hvor \VC sørger for en beholdning af Carlsberg 47 øl og \CERM skriver en hyldestsang til denne herlig øl. Et yndet foretagende til dette arrengement er at diskutere om markaten på Carlsberg 47 er blå eller grøn.

\subsection{En Kasse I En Festforening}
I lang tid har \TKET deltaget i \alkymias ølstafet. \TKETs medlemmer fandt dog hurtigt ud af, at løbe-delen af stafetten var den sværeste del. Dette førte til at \TKET i efteråret 2011 udfordrede kemikerne til en anderledes form for stafet: En ølstafet uden løb! For at kompensere for den reducerede mængde af fysisk aktivitet blev mængden der skulle drikkes sat op, hvilket endeligt gjorde det muligt at finde svar på spørgsmålet: Hvem kan drikke en kasse øl hurtigst! Derved startede \textit{En Kasse I En Festforening}.

I starten var konkurrencen simpel: To hold af seks personer dystede på hvor hurtigt de kunne tømme 30 øl. Dette var en stor success, hvor \TKET vandt 8 gange på de første 10 år, med den hurtigste tid på 1:12 sat i 2019\footnote{Dog fik \TKET en tid på 1:09 i 2021, hvor formatet var blevet tilpasset ift. corona-nedlukningen. Konkurrencen blev live-streamede, og atleterne drak ``1 mod 1" over et zoom-opkald. Efter alle havde drukket færdigt blev tiderne samlet til én fælles tid for hvert hold. Her vandt \TKET med den nævnte tid på 1:09 mod \alkymias tid på 1:41.}.
For at gøre konkurrencen mere tilskuervendelig\footnote{De lave tider var imponerende at overvære, men en konkurrence hvor holdenes introduktion var længere end selve dysten, blev dømt til at gå for hurtigt.}, valgte man i 2022 at ændre formatet. Holdene blev udvidet fra 6 personer til 8. og i stedet for, at alle måtte drikke på samme tid, gik man tilbage til oprindelsen af dysten, og lavede det om til en reel stafet hvor øllene skal drikkes sekventielt! Dette gav de mange tilskuere mere tid til at heppe på deres forening. Efter det nye format blev indført er den hurtigste tid sat i 2025 på 7:01. 

\subsection{MGP}
Det skete i de dage at en underforening så børnenes MGP på DR og tænkte: det er nogle gode og sjove melodier, men teksterne kunne blive bedre og sjovere! Det blev startskuddet til det første MGP i foråret 2014 i Staff Lounge. Fire sange stillede op første år og vinderen blev 'Allah, tro på os to'. Arrangementet blev en succes og over årene er det kun blevet større. I 2018 blev det muligt at stemme via SMS og i 2019 flyttede arrangementet i Matlab for at få plads til alle de trofaste deltagere og publikummer.

MGP starter altid med at se DR's MGP live i et auditorium (med passende drukregler) for at finde inspiration til næste års melodier, inden vi rykker videre til vores eget MGP. Nogle af de ikoniske sange gennem årene inkluderer titler som 'Bundebajer', 'Ti arme', 'Gnu er (det du lugter af)' og 'Hvis FU ku' bestemme'. 

\epigraph{Karina Sunds Nielsen,\\
SEKR 2013-2014}

\subsection{Æggepustning}
Påsken er om hjørnet og det betyder selvfølgelig, at der skal pyntes op på \TKET. Derfor inviterer PR til æggepustning, hvor der skal pustes og dekoreres æg. Som tak for fremstilling af en malet æg til ophængning på \TKET sørger PR for rigeligt med LFPandekager og hvid glögg.

%\paragraph{Kapsejlads}
%Engang i tidernes morgen besluttede nogle læger og tandlæger sig for at forstyrre ænderne i unisøen med en ølstafet til vands. Det var første gang \TKET vandt kapsejladsen. Siden har vi rent faktisk også fået lov til at deltage, og arrangementet er nu så stort, at det tiltrækker over 30.000 tilskuere. Dette gav visse problemer med at skaffe pladser, hvor man rent faktisk kunne se \TKETs suveræne sejr, hvorfor innovative hænger med hang til værktøj i ugerne op til Kapsejladsen konstruerede en prægtig flåde, så de kunne nyde \TKETs overlegenhed fra første parket midt på søen. Hængerflåden blev dog fjernet i XXX for i stedet at gøre plads til den mere inkluderende XXX camp i samarbejde med XXX, som først så dagens lys i XXX. 2026 blev året hvor vi for første gang i nyere tid deltog i Kapsejladsen på helt lige fod med de andre foreninger.

\subsection{TK-Open}
TK-open er \TKETs årlige brætspilsaften. Her er alle velkomne til en aften i MatLab, hvor diverse spil ligges frem, og baren er åben med drinks, øl og sodavand. Der bliver arrangeret turneringer i alt fra skak og neutrino til 4-på-stribe i samtlige dimensioner og vektorrally, med mere. Aftenens højdepunkt er dog for mange øørllhhludooo, hvor PR har lavet en smuk ludoplade, som dog hurtigt bliver dækket af øl og bræk, når 4 hold, af to personer på hver, kæmper om at snyde mest i ludo, og få deres 4 spillebrikker, som alle er øl, først i mål, alt imens \CERM og FORM agere kommentatorer og dommer, med gavmild uddeling af straf bundeøl.

\subsection{Skabsbundingstur}
Studerende bor ofte småt og bliver derfor nødt til at pladsoptimere deres hjem. Hertil yder hængerne på \TKET en speciel service: at teste hvor vidt et skab eller skabs-ækvivalent har en passende størrelse. Skabstesten kan udføres på mange måde men den mest præcise og mest elegante metode, som hængerne har fundet, er at have åbne drikkevarer parat, pakke så mange hængere som muligt ind i skabet, test om skabet kan lukkes, hvorefter samtlige hængere i skabet bunder de førnævnte drikkevarer. Dette foregår et par gange om året hvor hængerne drager ud i bybilledet og går fra lejlighed til lejlighed og tester hinandens skabe.

\subsection{Ølsmagning}
De sidste par år har \TKET afholdt en ølsmagning, hvor foreningen har inviteret et bryggeri forbi Matematisk Institut. Indtil videre har vi haft fornemt besøg af Mikkeller, Ebeltoft Gårdbryggeri og Ølsnedkeren. Efter ølsmagningen blev der afholdt Café Bagefter, hvilket altid er endt festligt. I 2025 blev der afholdt Café Bagefter i MatKant i samarbejde med Kalkulerbar, hvilket kun øgede den festlige stemning. 

\subsection{Majrevy og -fest}
Den anden halvårlige revy og tilhørende fest. Da der også sker ting på nationalt og internationalt plan i foråret, bliver der hurtigt indsamlet materiale til en ny revy. Til denne revy holder fOrm en tale til folket, og minder os om at \TKET vil bestå. Revyen afsluttes med ``Kom, maj, du søde, milde'' og efterfølgende fest.

\subsection{Lanciers}
Træk i dit stiveste puds og kom til en hyggelig aften med dans i Matlab, når Snuptaw spiller op til Les Lanciers. Her lærer vi at danse de forskellige runder i lanciers, og der er selvfølgelig også mulighed for at danse en omgang bukseløs can-can eller to.

\subsection{FU-Løb}
Den sædvanlige løbeseddelrute løbes igennem ved supersoniske hastigheder af alle FU’er, samt forsvarende mester. Undervejs løses opgaver på poster. Det er tilladt at fedte og på andre måder bestikke dommerne.

\subsection{Bundeservice}
Bundeservice er et velgørenhedsarrangement afholdt af en underforening, hvor at afholderne, og andre tørstige personer i nærheden af \TKET, bunder øl for velgørenhed. Til hver afholdelse af Bundeservice bliver der valgt et velgørende formål, f.eks. Røde Kors eller Kræftens Bekæmpelse, som der doneres til. En person kan så ringe ind til \TKET{}, og donere en drikkevare til en tørstig bunder, som så donere genstandens pris til det velgørende formål. 

\subsection{Kaløturen}
I sommerferien begiver en flok entusiastiske kammergængere sig til Kalø på en weekend campingtur. Der skal hygges, drikkes nogle drikkevarer, og måske svømmes med en bowlingkugle. Derudover besøges vagtbrien Brians gravsted ved Kalø Slotsruin, hvor der fortælles historier og synges sange til ære for ham.

\subsection{MatlaBotanisk Have}
I løbet af sommeren tages en tur i Botanisk Have, hvor hængere spiller ølbowling, spiser smørrebrød og drikker diverse drikkevarer i solens stråler. Skulle solen udeblive flyttes arrangementet typisk indendørs til Matlab.

\subsection{SØHYG}
Tidligere kaldet SØØL. Dette er rusernes første kontakt med \TKET. Tutorerne fra Mat/Fys-Tutorforeningen tager deres hold med nye studerende med ned til søen i universitetsparken, hvor \TKET byder på drikkevarer. Hvis vejret ikke arter sig, flyttes søen dog indenfor i Matlab.

\subsection{Rushyg}
Dette er en fest kun for de nye studerende og deres tutorer. Her kan russerne hygge sig med deres tutorer, og ligeledes forsøge at drikke deres tutorer fulde. \TKET giver per tradition russerne én øl.

\subsection{Pub-Quiz}
Pub-quiz er et arrangement arrangeret af hængere. Her skal der på hold quizzes i stort og småt og alt derimellem. Under og efter quizzen er der mulighed for at købe lækre drikkevarer samt sidde og hygge.

\subsection{En Kasse i \KASS}
Hvert år vælges en ny \KASS, og hvert år presser det samme spørgsmål sig derfor på. Kan der være en kass' i \KASS? For at svaret kan være ja, må \KASS indtage 30 øl uden at falde i søvn og uden at kast op.

\subsection{Generalforsamling}
Den tidligere BEST går af, og en ny BEST vælges på \textit{demokratisk} vis. FU'er vælges ved samme lejlighed blandt de nye studerende, og EFU'er bliver valgt blandt de 2. års studerende på Datalogi/IT.

%%%%%
\section{Tidligere afholdte arrangementer}

\subsection{Stiftelsesfest}
Festen har været afholdt fra \TKETs første år, og blev normalt afhold i starten af oktober -- hvor \TKET jo har fødselsdag -- og festens navn kan dateres helt tilbage til 1960. Stiftelsesfesten var derved den nye bestyrelses første fest. Her skulle den nye ceremonimester indsættes, og der var tradition for at høn til denne fest skulle vise sine evner indenfor alkoholindtagningens kunst. Her skulle \CERM igennem prøvelser involverende øl og sang, og alle var opfordret til at skåle med \CERM hele aftenen. Selvom festen ikke længere eksisterer i samme forstand som tidligere, har den nye Halloweenfest stadig medtaget visse dele af denne tradition i varierende omfang.

\subsection{Karnevals- og Kyndelmisse-fest}
Gennem tiden er der en håndfuld fester der er blevet slået sammen eller har ændret facon.

Ved Karnevalsfest fik de studerende lov at vise deres hemmelige identitet/tøjstil frem. Festen var en tradition der gik tilbage til 1963. Sidst i 1980’erne begyndte man at afholde festen sammen med andre naturvidenskablige festforeninger, hvilket i tidens løb har inkluderet @lkymia, Brocas Bodega, CHAOS, Fredagscaféen, MØF og Nanorama.

Kyndelmissefesten, også kendt som Kissemissefest, blev afholdt frem til 2020, men grundet det lille antal gæster blev denne fest slået sammen Karnevalsfesten.

Karnevalsfesten blev i midt 20'erne moderniseret i forbindelse med ændringerne til Stiftelsesfesten og idag afholdes Halloweenfesten i stedet.

\newpage

\section{PET-brev}
Engang i det herrens år 1998 fik den daværende næstformand en lys ide, og sendte
dette brev til PET:
\epigraph{\TKET, 22/11 1998}

\tkpicture[0.8\linewidth]{../img/pet_brev}

Ovenstående resulterede i et opkald til \TKET indledt med ordene: „Det er
kriminalpolitiet. Kan jeg snakke med Hr. Søren Sandmann?“ De ville vist bare lige
sikre sig, at det var en spøg!
